\subsubsection*{Duas facturas na mesma carteira}

Bob entregou a Alice dois sub-endereços, um para cada factura.
Ela deseja pagá-las na mesma transacção. Suponhamos que os valores
condizem exactamente com as entradas e a comissão, de modo que
há duas saídas e ambas se destinam a Bob. Como há-de ele atribuir
cada saída à factura certa se a chave privada de vista é a mesma?

Conservemos as chaves didácticas de Bob, \(k_B^v=5\) e \(k_B^s=7\),
da secção~\ref{sec:dois-destinos-varrimento}. À primeira factura
coube o índice \((0,1)\) da carteira, à segunda o índice \((0,2)\).
Escrevamos \(S_j=K_B^{s,(0,j)}\) e \(V_j=K_B^{v,(0,j)}\), com
\(j\in\{1,2\}\). Bob pode preparar de antemão as duas linhas
\[
                 S_1\longmapsto(0,1),\qquad
                 S_2\longmapsto(0,2).
\]
Ambas se obtêm da chave pública principal \(7G\). Para a primeira,
a entrada de \(\mathcal H_n\) contém, por esta ordem, os oito bytes
de \(T_{\rm sub}\), os 32 bytes de \(5\) e os oito bytes dos
índices, estes últimos sob a forma
\(\texttt{00 00 00 00}\mathbin\|\texttt{01 00 00 00}\).
Para a segunda linha, o \texttt{01} passa a \texttt{02}. Em cada
caso, o escalar obtido multiplica \(G\), e o ponto resultante
soma-se a \(7G\). Eis as duas entradas da tabela e as suas
companheiras de vista, \(V_j=5S_j\). O quadro dá os primeiros
oito bytes de cada codificação, para podermos reconhecer os
pontos mais adiante. As contas usam sempre os 32 bytes.
\[
\begin{array}{c|ll}
j&S_j&V_j\\ \hline
1&\texttt{89bf16f6674c1e80}&\texttt{2647edd726a81314}\\
2&\texttt{12b4613d46598f36}&\texttt{436768693cf4d097}
\end{array}
\]

Alice coloca a primeira factura na posição \(t=0\) da transacção e
a segunda na posição \(t=1\). Os índices \(j\) pertencem ao livro
de facturas de Bob. Os índices \(t\) contam as saídas que Alice
acaba de ordenar. A relação \(j=t+1\) é uma coincidência deste
pagamento. Outro remetente pode usar um só destes sub-endereços,
numa posição inteiramente diferente, sem saber sequer o número da
factura.

Os dois destinatários textuais são distintos. Além da chave
principal \(R\), Alice publica, pois, uma chave adicional por
saída. Para o exemplo, tomemos \(R=3G\), \(r_0=17\) e \(r_1=23\).
A Alice chegaram apenas os pares de pontos \((S_1,V_1)\) e
\((S_2,V_2)\). A partir deles calcula
\[
\begin{array}{c|cc}
 t&R_t&\mathcal D_{B,t}\\ \hline
 0&17S_1&8\cdot17V_1\\
 1&23S_2&8\cdot23V_2 .
\end{array}
\]
A chave de vista privada de Bob aparece de uma só vez nas duas
contas. Na primeira posição, \(8\cdot17V_1=8\cdot17(5S_1)
=40R_0\). Na segunda, \(8\cdot23V_2=40R_1\). Eis o que dão as
codificações e as etiquetas, com a mesma abreviatura de oito bytes
para os pontos. As etiquetas figuram por inteiro.
\[
\begin{array}{c|lll}
 t & R_t & \mathcal D_{B,t} & \text{etiqueta}\\ \hline
 0 & \texttt{42a01e4a01c93557} & \texttt{db0bb99b5ed23ea0} & \texttt{9d}\\
 1 & \texttt{b1d8f4d393f3eee5} & \texttt{461a4535027b57e3} & \texttt{4e}
\end{array}
\]

As chaves que Bob encontra nas saídas começam por
\(\texttt{557e9f5510f08d81}\) e \(\texttt{481dc3f4fcd3afd3}\).
Começa o varrimento pela chave principal \(R=3G\). Já conhecemos
a sua derivação para Bob, \(120G\), e as etiquetas que dá para os
dois índices, \texttt{92} e \texttt{45}. A primeira saída pede
\texttt{9d} e a segunda \texttt{4e}. Bob passa em ambas à chave
adicional da posição respectiva. Com \(40R_0\) e \(t=0\) encontra
\texttt{9d}. Com \(40R_1\) e \(t=1\) encontra \texttt{4e}.
Chegou a hora de decidir a que factura pertence cada pagamento.

Na primeira saída a carteira calcula
\(h_0=\mathcal H_n(\mathcal D_{B,0},\texttt{00})\) e subtrai
\(h_0G\) à chave recebida. A soma que Alice fizera era
\(K_0^o=h_0G+S_1\), pelo que a subtracção restitui o ponto que
começa por \texttt{89bf16f6674c1e80}. Na segunda, o escalar
\(h_1=\mathcal H_n(\mathcal D_{B,1},\texttt{01})\) aplicado à
mesma operação restitui o ponto que começa por
\texttt{12b4613d46598f36}. Bob procura os pontos completos na
tabela. Lê na primeira linha \((0,1)\), na segunda \((0,2)\), e
pode dar as duas facturas por pagas. O número de factura atravessou
toda a operação no registo privado de Bob. Na transacção viajaram
as chaves que lhe permitiram voltar a encontrá-lo.

Suponhamos agora que o registo de Bob contém dez mil facturas.
A preparação da tabela compreende dez mil derivações de
sub-endereços. Depois de preparada, o trabalho para classificar
estas duas saídas conserva a mesma forma. A carteira experimenta
a chave de transacção pertinente e o índice da saída, recupera o
ponto de gasto e procura-o entre as chaves da tabela. Uma tabela
de dispersão bem dimensionada permite a cada procura um número
esperado de operações independente da quantidade de facturas.
Se Bob restaurar a carteira com uma tabela que abranja somente a
primeira factura, a segunda procura chegará ao ponto correcto,
mas faltará a linha que lhe dá o nome. Reconstituir essa linha a
partir de \((0,2)\) e voltar a examinar a saída bastará para a
reconhecer. A chave de vista determina a conta comum. A extensão
da tabela determina quais dos seus sub-endereços a carteira pode
identificar de imediato.
